% © 2026 Ing. David Jaroš — CC BY-NC-ND 4.0
%
% This work is licensed under a Creative Commons Attribution-NonCommercial-NoDerivatives
% 4.0 International License (CC BY-NC-ND 4.0).
%
% License History: Earlier drafts (up to v0.3) were released under CC BY 4.0.
% From v0.4 onward, all material is released under CC BY-NC-ND 4.0 to protect
% the integrity of the theoretical work during ongoing academic development.
%
% See LICENSE.md for full license text.
%
% File: research_tracks/T3_ALPHA/nu2_correction_derivation.tex
%
% Purpose: Track the derivation attempt of the modular elliptic correction
%   +nu2(p)/4 in B(p) from the UBT action S[Theta] via Z2-twisted sectors.

% UBT-AI-PROVENANCE-BEGIN
% schema: ubt-ai-provenance/v1
% tier: C_working
% ai_assistance: disclosed
% human_review: risk-based
% editorial_responsibility: Ing. David Jaroš
% policy: ../../AI_PROVENANCE.md
% notice: Working material; exhaustive human review is not claimed.
% UBT-AI-PROVENANCE-END

\ifdefined\INCLUDEMODE
\else
  \documentclass[12pt,a4paper]{article}
  \usepackage[a4paper, margin=2.5cm]{geometry}
  \usepackage{amsmath, amssymb}
  \usepackage{hyperref}
  \usepackage{xcolor}
  \usepackage{tcolorbox}
  \tcbuselibrary{skins,breakable}

  \newtcolorbox{statusbox}[1][]{colback=gray!8!white,colframe=black!60,
    arc=3pt,boxrule=1pt,fonttitle=\bfseries,title=Status Summary,#1}

  \title{\textbf{Derivation Attempt of the \(\nu_2(p)/4\) Correction from \(S[\Theta]\)}}
  \author{Ing.\ David Jaro\v{s}}
  \date{2026-05-11}
  \begin{document}
  \maketitle
\fi

\section{Target statement}

The target is to derive, from UBT action data (without $\alpha$ input),
\begin{equation}
  B(p) = \frac{p+1}{3} + \frac{\nu_2(p)}{4},
  \qquad
  \nu_2(p) = 1 + \left(\frac{-4}{p}\right),
\end{equation}
with $\nu_2(p)=2$ for $p\equiv 1 \!\!\pmod 4$ and $\nu_2(p)=0$ for
$p\equiv 3 \!\!\pmod 4$.

\section{Working mechanism under evaluation}

We evaluate the hypothesis that the involution
\(\tau:\Theta\mapsto\Theta^\dagger\) defines a physically active
\(\mathbb{Z}_2\)-twisted sector in the winding background \(\Theta_n\), and that
the one-loop determinant on the \(\tau\)-fixed fluctuation submanifold
\(\delta\Theta=\delta\Theta^\dagger\) contributes an orbifold term matching the
Riemann--Hurwitz elliptic correction scale \(1/4\).

\section{Current technical state}

\begin{itemize}
  \item Arithmetic side (\(\nu_2\) formula, modular interpretation): established.
  \item UBT side (explicit determinant factorisation into untwisted/twisted
    sectors with coefficient fixed to \(+\nu_2/4\)): not yet completed.
  \item Non-circularity check R1--R5: currently satisfied at formulation level
    (no \(\alpha_{\mathrm{exp}}\), no fixed \(p=137\) input in the derivation ansatz).
\end{itemize}

\section{Immediate bridge conditions}

To promote from conjectural structure to derivation, the following must be shown:
\begin{enumerate}
  \item Explicit decomposition of one-loop operator
    \(\Delta_\Theta = \Delta_{\mathrm{untw}} \oplus \Delta_{\mathrm{tw}}\)
    induced by \(\tau\).
  \item Modular counting of twisted fixed classes matched to
    \(\nu_2(p)\) in the \(\Gamma_0(p)\) sector.
  \item Normalisation of the twisted contribution to \(B\) yielding exactly
    \(+\nu_2/4\), not a free prefactor.
\end{enumerate}

\begin{statusbox}
\textbf{Claim ``\(+\nu_2(p)/4\) is derived from \(S[\Theta]\)''}: \textbf{[OPEN]}.

\textbf{What is currently established}: modular arithmetic identity and physically
motivated \(\mathbb{Z}_2\)-twist mapping [MC].

\textbf{Blocking gap}: first-principles determinant normalisation in UBT action
language is still missing.
\end{statusbox}

\ifdefined\INCLUDEMODE
\else
  \end{document}
\fi
